\begin{abstract}
Hybrid ReRAM crossbars keep the most important weights in reliable single-level cells (SLC) and the rest
in dense multi-level cells (MLC). HyFlexPIM applies this to transformers by factorising each weight
matrix with SVD and protecting the components with the largest loss gradients. We evaluate this scheme on
BERT-Base and DistilBERT (SST-2) with a bit-accurate crossbar simulator that models errors cell by cell:
programming and read noise, level-dependent noise, stuck cells, a fixed ADC with bit-serial inputs and IR
drop. A hybrid layout is necessary, since all-MLC storage falls to 52--60\% accuracy at 7.5\% programming noise,
but the choice of components matters as much as their number. Ranking components by singular value
reaches the same accuracy as a Fisher-information estimate of each component's expected damage, while a
gradient score loses 2--12 points at high noise; we explain this through the per-column quantisation of the
crossbar. With magnitude ranking, 5--10\% of components in SLC match an all-SLC chip at 5\% noise.
Stuck cells are the one effect placement cannot handle. Because the components of an SVD layer can be
stored in any order, we map them to physical slots around a tested chip's stuck cells by solving an
assignment problem, at no hardware cost; this raises BERT-Base at 0.5\% stuck cells from 76.6\% to
85.6\%. Fine-tuning on simulated chips with stuck cells adds to this mapping. A tier-aware variant, which
trains each component under the errors of the tier (SLC or MLC) it will be stored in, lifts BERT-Base at 1\% stuck cells from 65.1\% to 82.9\%,
so that the hybrid chip matches or exceeds an all-SLC chip with mapping up to 1\% scattered stuck cells,
using 34\% fewer cells. The placement and mapping results also hold
on MNLI.
\end{abstract}

\begin{IEEEkeywords}
ReRAM, processing-in-memory, multi-level cell, stuck-at faults, SVD, transformer, fault-aware mapping, noise-aware training
\end{IEEEkeywords}
